% GENERATED FILE. Edit canonical lessons, then run npm run generate:books.
% canonical-source-hash: fnv1a64:c67435a1fb87662e
% canonical-lessons: RU-C372-veroyatno, RU-C372-veroyatnost, RU-C372-maloveroyatnyy, RU-C372-shans, RU-C372-risk

\chapter{Likely or unlikely}
\label{ch:ru-likely-or-unlikely}
\hlchaptermodality{372}

\begin{chapteropening}
\textbf{By the end of this chapter:} \emph{I can say that something is probable or unlikely, talk about probability, and speak about a chance and a risk.}

Complete the last lesson of chapter 372: I can say that something is probable or unlikely, talk about probability, and speak about a chance and a risk.
\end{chapteropening}

\section[\texorpdfstring{\ru{вероятно} (veroyátno)}{вероятно (veroyátno)}]{\ru{вероятно} (veroyátno) — probably, likely}

\label{lesson:RU-C372-veroyatno}



\begin{quote}
Before the new one: say the Russian for let us say, then the Russian for a supposition.
\end{quote}

\subsection*{You'll want to know: \ru{вероятно}}
\textbf{\ru{вероятно}} (\emph{veroyátno}) — \textquotedblleft{}probably, likely\textquotedblright{}.

\begin{quote}
\ru{Он}, \ru{вероятно}, \ru{опоздает}.
\end{quote}

\begin{quote}
\emph{On, veroyátno, opozdáyet.}
\end{quote}

\begin{quote}
He will probably be late.
\end{quote}

More bookish than \emph{\ru{наверное}}; set off by commas inside the sentence.

\begin{cousinweb}
\textbf{\ru{вероятно}} is built on \emph{\ru{вероятный}} \textquotedblleft{}likely\textquotedblright{}, which is usually explained as \emph{\ru{вера}} \textquotedblleft{}belief\textquotedblright{} with the old verb \emph{\ru{ять}} \textquotedblleft{}to take\textquotedblright{}: what may be taken on belief. \emph{\ru{Вера}} is inherited from Proto-Slavic \emph{*věra} (EtymDB).
\end{cousinweb}

\begin{grammarlens}[title={What you've built}]
One more word for telling, explaining and describing at B1.
\end{grammarlens}

\subsection*{Your turn}
\begin{itemize}
\raggedright
  \item \emph{Say it:} \emph{veroyátno}
  \item \emph{Say it:} \emph{veroyátno}, once more
  \item \emph{Say it:} a sentence of your own with \emph{veroyátno}
  \item \emph{Recall:} say the Russian for let us say, then the Russian for a supposition, then say \emph{veroyátno} again
\end{itemize}

\begin{tcolorbox}[breakable,colback=teal!4,colframe=teal!35!black,title={Before you move on}]
What does \ru{вероятно} mean? (\textquotedblleft{}Probably, likely\textquotedblright{}.) Say it once more.
\end{tcolorbox}
\section[\texorpdfstring{\ru{вероятность} (veroyátnost)}{вероятность (veroyátnost)}]{\ru{вероятность} (veroyátnost) — probability, likelihood}

\label{lesson:RU-C372-veroyatnost}



\begin{quote}
Before the new one: say the Russian for a supposition, then the Russian for probably.
\end{quote}

\subsection*{You'll want to know: \ru{вероятность}}
\textbf{\ru{вероятность}} (\emph{veroyátnost}) — \textquotedblleft{}probability, likelihood\textquotedblright{}.

\begin{quote}
\ru{Вероятность} \ru{дождя} — \ru{сорок} \ru{процентов}.
\end{quote}

\begin{quote}
\emph{Veroyátnost dozhdyá — sórok protséntov.}
\end{quote}

\begin{quote}
The chance of rain is forty per cent.
\end{quote}

Feminine. \emph{\ru{По} \ru{всей} \ru{вероятности}} — in all likelihood.

\begin{cousinweb}
\textbf{\ru{вероятность}} is built on \emph{\ru{вероятный}} \textquotedblleft{}likely\textquotedblright{} with the abstract suffix \emph{-\ru{ость}}; the root is \emph{\ru{вера}} \textquotedblleft{}belief\textquotedblright{}, Proto-Slavic \emph{*věra} (EtymDB).
\end{cousinweb}

\begin{grammarlens}[title={What you've built}]
One more word for telling, explaining and describing at B1.
\end{grammarlens}

\subsection*{Your turn}
\begin{itemize}
\raggedright
  \item \emph{Say it:} \emph{veroyátnost}
  \item \emph{Say it:} \emph{veroyátnost}, once more
  \item \emph{Say it:} a sentence of your own with \emph{veroyátnost}
  \item \emph{Recall:} say the Russian for a supposition, then the Russian for probably, then say \emph{veroyátnost} again
\end{itemize}

\begin{tcolorbox}[breakable,colback=teal!4,colframe=teal!35!black,title={Before you move on}]
What does \ru{вероятность} mean? (\textquotedblleft{}Probability, likelihood\textquotedblright{}.) Say it once more.
\end{tcolorbox}
\section[\texorpdfstring{\ru{маловероятный} (maloveroyátnyy)}{маловероятный (maloveroyátnyy)}]{\ru{маловероятный} (maloveroyátnyy) — unlikely}

\label{lesson:RU-C372-maloveroyatnyy}



\begin{quote}
Before the new one: say the Russian for probably, then the Russian for probability.
\end{quote}

\subsection*{You'll want to know: \ru{маловероятный}}
\textbf{\ru{маловероятный}} (\emph{maloveroyátnyy}) — \textquotedblleft{}unlikely\textquotedblright{}.

\begin{quote}
\ru{Это} \ru{маловероятно}, \ru{но} \ru{возможно}.
\end{quote}

\begin{quote}
\emph{Éto maloveroyátno, no vozmózhno.}
\end{quote}

\begin{quote}
It is unlikely, but possible.
\end{quote}

Built from \emph{\ru{мало}}, little, and \emph{\ru{вероятный}}, likely. \emph{\ru{Маловероятно}} — it is unlikely.

\begin{cousinweb}
\textbf{\ru{маловероятный}} is built from \emph{\ru{мало}} \textquotedblleft{}little\textquotedblright{}, from Proto-Slavic \emph{*malŭ} \textquotedblleft{}small\textquotedblright{} (EtymDB), and \emph{\ru{вероятный}} \textquotedblleft{}likely\textquotedblright{}, built on \emph{\ru{вера}} \textquotedblleft{}belief\textquotedblright{}.
\end{cousinweb}

\begin{grammarlens}[title={What you've built}]
One more word for telling, explaining and describing at B1.
\end{grammarlens}

\subsection*{Your turn}
\begin{itemize}
\raggedright
  \item \emph{Say it:} \emph{maloveroyátnyy}
  \item \emph{Say it:} \emph{maloveroyátnyy}, once more
  \item \emph{Say it:} a sentence of your own with \emph{maloveroyátnyy}
  \item \emph{Recall:} say the Russian for probably, then the Russian for probability, then say \emph{maloveroyátnyy} again
\end{itemize}

\begin{tcolorbox}[breakable,colback=teal!4,colframe=teal!35!black,title={Before you move on}]
What does \ru{маловероятный} mean? (\textquotedblleft{}Unlikely\textquotedblright{}.) Say it once more.
\end{tcolorbox}
\section[\texorpdfstring{\ru{шанс} (shans)}{шанс (shans)}]{\ru{шанс} (shans) — a chance}

\label{lesson:RU-C372-shans}



\begin{quote}
Before the new one: say the Russian for probability, then the Russian for unlikely.
\end{quote}

\subsection*{You'll want to know: \ru{шанс}}
\textbf{\ru{шанс}} (\emph{shans}) — \textquotedblleft{}a chance\textquotedblright{}.

\begin{quote}
\ru{У} \ru{тебя} \ru{есть} \ru{шанс} \ru{получить} \ru{эту} \ru{работу}.
\end{quote}

\begin{quote}
\emph{U tebyá yest shans poluchít étu rabótu.}
\end{quote}

\begin{quote}
You have a chance of getting this job.
\end{quote}

Masculine. \emph{\ru{Шансы}} — the odds: \emph{\ru{у} \ru{него} \ru{мало} \ru{шансов}}.

\begin{cousinweb}
\textbf{\ru{шанс}} was borrowed from French \emph{chance} \textquotedblleft{}luck, chance\textquotedblright{}, from Vulgar Latin \emph{cadentia} \textquotedblleft{}a falling\textquotedblright{}, from Latin \emph{cadere} \textquotedblleft{}to fall, to happen\textquotedblright{} (EtymDB): how the dice fall.
\end{cousinweb}

\begin{grammarlens}[title={What you've built}]
One more word for telling, explaining and describing at B1.
\end{grammarlens}

\subsection*{Your turn}
\begin{itemize}
\raggedright
  \item \emph{Say it:} \emph{shans}
  \item \emph{Say it:} \emph{shans}, once more
  \item \emph{Say it:} a sentence of your own with \emph{shans}
  \item \emph{Recall:} say the Russian for probability, then the Russian for unlikely, then say \emph{shans} again
\end{itemize}

\begin{tcolorbox}[breakable,colback=teal!4,colframe=teal!35!black,title={Before you move on}]
What does \ru{шанс} mean? (\textquotedblleft{}A chance\textquotedblright{}.) Say it once more.
\end{tcolorbox}
\section[\texorpdfstring{\ru{риск} (risk)}{риск (risk)}]{\ru{риск} (risk) — a risk}

\label{lesson:RU-C372-risk}



\begin{quote}
Before the new one: say the Russian for unlikely, then the Russian for a chance.
\end{quote}

\subsection*{You'll want to know: \ru{риск}}
\textbf{\ru{риск}} (\emph{risk}) — \textquotedblleft{}a risk\textquotedblright{}.

\begin{quote}
\ru{Это} \ru{большой} \ru{риск}.
\end{quote}

\begin{quote}
\emph{Éto bolshóy risk.}
\end{quote}

\begin{quote}
That is a big risk.
\end{quote}

Masculine. \emph{\ru{На} \ru{свой} \ru{риск}} — at one’s own risk.

\begin{cousinweb}
\textbf{\ru{риск}} is usually traced, through French \emph{risque}, to Italian \emph{risco}; the further origin of the Italian word is not settled.
\end{cousinweb}

\begin{grammarlens}[title={What you've built}]
One more word for telling, explaining and describing at B1.
\end{grammarlens}

\subsection*{Your turn}
\begin{itemize}
\raggedright
  \item \emph{Say it:} \emph{risk}
  \item \emph{Say it:} \emph{risk}, once more
  \item \emph{Say it:} a sentence of your own with \emph{risk}
  \item \emph{Recall:} say the Russian for unlikely, then the Russian for a chance, then say \emph{risk} again
\end{itemize}

\begin{tcolorbox}[breakable,colback=teal!4,colframe=teal!35!black,title={Before you move on}]
What does \ru{риск} mean? (\textquotedblleft{}A risk\textquotedblright{}.) Say it once more.
\end{tcolorbox}
